%!TEX TS-program = xelatex
%!TEX encoding = UTF-8
%!BIB TS-program = bibtex
%!TEX spellcheck = en_US

\documentclass[11pt,a4paper]{article}
\usepackage{amsmath}
\usepackage{amsfonts}
\usepackage{amssymb}
\usepackage{graphicx}
\graphicspath{{images/}}
\usepackage{enumerate}
\usepackage{braket}

\title{Some comments and figures on Homework 1}		% Title
\author{}								% Author
\date{\today}											% Date

\begin{document}
\maketitle
\section{Problem 1}
\begin{enumerate}[(a)]
	\item See Fig. \ref{fig:pro1a} for reference.
	      \begin{figure}[h]
		      \centering
		      \includegraphics[width=0.7\linewidth]{pro_1_a.pdf}
		      \caption{Parameters: $k_B = \gamma = 1$.}
		      \label{fig:pro1a}
	      \end{figure}
	\item See Fig. \ref{fig:pro1b} for reference.
	      \begin{figure}[h]
		      \centering
		      \includegraphics[width=0.7\linewidth]{pro_1_b.pdf}
		      \caption{Parameters: $k_B = \gamma = 1$.}
		      \label{fig:pro1b}
	      \end{figure}
	\item See Fig. \ref{fig:pro1c} for reference.
	      \begin{figure}[h]
		      \centering
		      \includegraphics[width=0.7\linewidth]{pro_1_c.pdf}
		      \caption{Parameters: $k_B = \gamma = 1$, and $\frac{1}{4\hbar} \big( \frac{m}{\gamma} \big) ^{\frac{1}{2}}= 1$.}
		      \label{fig:pro1c}
	      \end{figure}
\end{enumerate}

\newpage
\section{Problem 2}
\begin{enumerate}[(a)]
	\setcounter{enumi}{2}
	\item To make it simple, I assume $k_B T = \varepsilon$.
	      And let the analytic result is
	      \begin{equation}
		      \braket{x} = \frac{4\lambda k_B T}{3 \alpha k_B T - 4 \gamma^2} =
		      \frac{4\lambda \varepsilon}{3 \alpha \varepsilon - 4 \gamma^2},
	      \end{equation}
	      and mean-field result is
	      \begin{equation}
		      \braket{x}_{\text{mean-field}} =
		      \frac{1}{3} \frac{4\lambda k_B T}{3 \alpha k_B T - 4 \gamma^2} =
		      \frac{1}{3} \frac{4\lambda \varepsilon}{3 \alpha \varepsilon - 4 \gamma^2}.
	      \end{equation}

	      Perform molecular dynamics simulation,
	      when $\braket{x}$ is very small,
	      we can see analytic result is closer to the simulation result,
	      but mean-field has noticeable deviation from simulation;
	      when $\braket{x}$ is large, mean-filed result is closer.
	      See Fig. \ref{fig:pro2c2} for reference.
	      \begin{figure}[h]
		      \centering
		      \includegraphics[width=0.7\linewidth]{pro_2_c_2.pdf}
		      \caption{Simulation result compare to analytic and mean-filed results.}
		      \label{fig:pro2c2}
	      \end{figure}
	      Here the initial $x$ position is from $0$ to $0.5$, with $2^{14}$ time steps.
\end{enumerate}


\end{document}
